\documentclass[../../../include/open-logic-section]{subfiles}

\begin{document}

\olfileid{sth}{replacement}{ref}
\olsection{替换与反射}

我们借助 \stagesinex{} 为替换公理提供辩护的最后一次尝试，始于一个深刻而优美的结果：\footnote{提醒：所有公式均可带参数（除非另有明确说明）。}
%In this, I will use overlining, such as $x_1, \ldots, x_n$, to
%abbreviate “$x_1, \ldots, x_n$”:


\begin{thm}[反射模式]\ollabel{reflectionschema}
对任意公式 $\phi$：
\[
\forall \alpha \exists \beta > \alpha (\forall x_1 \ldots, x_n \in
V_\beta)(\phi(x_1, \ldots, x_n) \liff \phi^{V_\beta}(x_1, \ldots, x_n))
\]
\end{thm}


\noindent
如第 \olref[sth][replacement][strength]{formularelativization} 节所述，$\phi^{V_\beta}$ 是将 $\phi$ 中的所有量词限制到集合~$V_\beta$ 所得的结果。因此，直观上，反射说的是：如果 $\phi$ 在整个层级中为真，那么 $\phi$ 在该层级任意多个\emph{初始段}中亦为真。

\citet{Montague1961} 和 \citet{Levy1960} 证明了，（适当表述下的）替换公理与反射在~$\Z$ 之上等价，因此添加其中任何一个皆可得到~$\ZF$。（我们将在第 \olref[sth][replacement][refproofs]{sec} 节证明这些结果。）鉴于这种等价性，人们或许能如下地借助 \stagesinex{} 为反射与替换公理提供辩护：根据 \stagesinex{}，层级应当极高；事实上高到，我们对其所能作出的任何断言都不足以限制其高度。我们可以将此理解为：如果任一语句~$\phi$ 在整个层级中为真，那么它在层级任意多个初始段中亦为真。而这正是反射。

同样，这似乎是一种颇具前景的尝试，旨在为替换公理提供内在辩护。但关于这一点可讨论的内容实在太多。现在，你必须自行判断这种尝试是否成功。\footnote{不过，你可能想继续阅读 \citet[95--100]{Incurvati2020}。}

\end{document}